\makeatletter
\newcommand{\oset}[3][0ex]{
	\mathrel{\mathop{#3}\limits^{\vbox to#1{\kern-2\ex@\hbox{$\scriptstyle#2$}}}}
}
\makeatother

% Expressions -------------------------------------------------------------------------------------------------------

\newcommand{\expr}[1]{%
    \Colorbox{backcolor}{\ensuremath{\mathtt{ #1 }}}%
}

% MMT -----------------------------------------------------------------------------------------------------------------
% judg, mmt

\newcommand{\judgmentscolor}{RedViolet}
\newcommand{\mmtcolor}{Turquoise}

\newcommand{\oma}[2]{\ensuremath{#1\left(#2\right)}}
\newcommand{\ombind}[3]{\ensuremath{#1\left[#2\right]\left(#3\right)}}

\newcommand{\judg}[2]{#1 \mylink{\judgmentscolor}{judg}{\vdash} #2}
\newcommand{\gjudg}[1]{\judg\Gamma{#1}}
\newcommand{\tjudg}[2]{\judg{\Theta;#1}{#2}}
\newcommand{\tgjudg}[1]{\tjudg{\Gamma}{#1}}
\newcommand{\infertype}[2]{ #1 \mylink{\judgmentscolor}{judg}{\Rightarrow} #2 }
\newcommand{\hastype}[2]{ #1 \mylink{\judgmentscolor}{judg}{\Leftarrow} #2 }
\newcommand{\judgequt}[2]{ #1 \mylink{\judgmentscolor}{judg}{\equiv} #2}
\newcommand{\judgeq}[3]{\judgequt{#1}{#2}\;\mylink{\judgmentscolor}{judg}{:}#3}
\newcommand{\judgeqc}[2]{ #1 \mylink{\judgmentscolor}{judg}{\rewrites} #2}
\newcommand{\subtp}[2]{ #1 \mylink{\judgmentscolor}{judg}{<:} #2}
\newcommand{\wfctx}[1]{\vdash #1\;\mylink{\judgmentscolor}{judg}{\mathtt{ctx}}}
\newcommand{\gwfctx}{\wfctx\Gamma}
\newcommand{\judginh}[1]{#1\;\mylink{\judgmentscolor}{judg}{\mathtt{inh}}}
\newcommand{\judguniv}[1]{#1\;\mylink{\judgmentscolor}{judg}{\mathtt{univ}}}
\newcommand{\judgmatchinf}[2]{ #1 \mylink{\judgmentscolor}{judgmatching}{\oset[-.45ex]\Rightarrow\equiv} #2 }

\newcommand\incl[1]{\mylink{\mmtcolor}{mmt}{\mathtt{include}}\ #1}
\newcommand\viewto[2]{#1 \mylink{\mmtcolor}{mmt}{\to} #2}
\newcommand\mmttp[1]{\mylink{\mmtcolor}{mmt}{\;:\;}#1}
\newcommand\mmtdf[1]{\mylink{\mmtcolor}{mmt}{\;=\;}#1}
\newcommand\mmtass[1]{\mylink{\mmtcolor}{mmt}{\;\mapsto\;}#1}


% LF ---------------------------------------------------------------------------------------------------------------------
%lfpi, jat,max

\newcommand{\lfcolor}{ForestGreen}
\newcommand\dedcolor{BurntOrange}

\newcommand{\umax}[1]{\mylink{\lfcolor}{max}{\max\left\{ \vphantom{#1} \right. } #1 \mylink{\lfcolor}{max}{\left. \vphantom{#1} \right\} }}

\newcommand{\type}{\mylink{\lfcolor}{lfpi}{\ensuremath{\mathtt{type}}}\xspace}
\newcommand{\kind}{\mylink{\lfcolor}{lfpi}{\ensuremath{\mathtt{kind}}}\xspace}
\newcommand\lfpisym{\mylink{\lfcolor}{lfpi}{\ensuremath{\prod}}\xspace}
\DeclareMathOperator*\lfpisymm{\mylink{\lfcolor}{lfpi}{\ensuremath{\prod}}}
\newcommand\lfpi[2]{\lfpisymm_{#1}#2}
\newcommand\lflambdasym{\mylink{\lfcolor}{lfpi}{\ensuremath{\lambda}}}
\newcommand\lflambda[2]{\lflambdasym_{#1}.\;#2}
\newcommand\lfarrowsym{\mylink{\lfcolor}{lfpi}{\ensuremath{\to}}}
\newcommand\lfarrow[2]{#1\; \lfarrowsym\; #2}

\newcommand\ded{\mylink{\dedcolor}{jat}{\textsc{\smaller DED}\;}\xspace}

% Sigma ------------------------------------------------------------------------------------------------------------------------------
% sigmatp

\newcommand\sigmacolor{Brown}

\newcommand{\sigmasymnolink}{\ensuremath{\sum}\xspace}
\newcommand{\sigmasym}{\hyperlink{sigmatp}{\textcolor{\sigmacolor}{\sigmasymnolink}}\xspace}
\DeclareMathOperator*{\sigmasymm}{\hyperlink{sigmatp}{\textcolor{\sigmacolor}{\sigmasymnolink}}\xspace}
\newcommand{\sigmatp}[2]{\sigmasymm_{#1}\;#2}
\newcommand{\pair}[2]{\mylink{\sigmacolor}{sigmatp}{\left\langle \vphantom{#1 #2} \right.} #1\hyperlink{sigmatp}{\textcolor{\sigmacolor}{\;,\;}} #2 \mylink{\sigmacolor}{sigmatp}{\left. \vphantom{#1 #2}\right\rangle}}
\newcommand{\projlsym}{\hyperlink{sigmatp}{\textcolor{\sigmacolor}{\ensuremath{\pi_\ell}}}\xspace}
\newcommand{\projl}[1]{\mylink{\sigmacolor}{sigmatp}{\projlsym\left( \vphantom{#1} \right.}#1\mylink{\sigmacolor}{sigmatp}{\left. \vphantom{#1}\right)}}
\newcommand{\projrsym}{\hyperlink{sigmatp}{\textcolor{\sigmacolor}{\ensuremath{\pi_r}}}\xspace}
\newcommand{\projr}[1]{\mylink{\sigmacolor}{sigmatp}{\projrsym\left( \vphantom{#1} \right.}#1\mylink{\sigmacolor}{sigmatp}{\left. \vphantom{#1}\right)}}
\newcommand{\simpleprod}[2]{#1 \hyperlink{sigmatp}{\textcolor{\sigmacolor}{\times}} #2}

% Coproducts ----------------------------------------------------------------------------------------------------------------------
% coprod, funcadd

\newcommand\coprodcolor{NavyBlue}

\newcommand{\coprodsym}{\mylink{\coprodcolor}{coprod}{\ensuremath{\bigoplus}}\xspace}
\renewcommand{\coprod}[2]{#1 \coprodsym #2}
\newcommand{\injlsym}{\mylink{\coprodcolor}{coprod}{\ensuremath{\hookrightarrow_\ell}}\xspace}
\newcommand{\injrsym}{\mylink{\coprodcolor}{coprod}{\ensuremath{\prescript{}{r}\hookleftarrow}}\xspace}
\newcommand{\injl}[2]{#1 \injlsym #2}
\newcommand{\injr}[2]{#2 \injrsym #1}
\newcommand{\cmatchsymnolink}{\ensuremath{\textbf{match}}\xspace}
\newcommand\cmatchsym{\mylink{\coprodcolor}{coprod}{\cmatchsymnolink}\xspace}
\newcommand{\cmatch}[5]{\cmatchsym_{#1}\mylink{\coprodcolor}{coprod}{\left\{ \vphantom{#2 #3 #4 } \right. } #2 \mylink{\coprodcolor}{coprod}{\Longrightarrow}_{#5}\; #3 \; \mylink{\coprodcolor}{coprod}{\left. \vphantom{#2 #3 #4} \middle|\; \right.} #4 \mylink{\coprodcolor}{coprod}{\left. \vphantom{#2 #3 #4} \right\} } }
\newcommand{\funcadd}[2]{#1 \mylink{\coprodcolor}{funcadd}{\oplus} #2}

% Finite Types-------------------------------------------------------------------------------------------------------------------------
% emptytype,unit,enum

\newcommand{\finitetypecolor}{Magenta}

\newcommand{\emptytype}{\ensuremath{\mylink{\finitetypecolor}{emptytype}{\emptyset}}\xspace}
\newcommand{\unitnolink}{\ensuremath{\textsc{Unit}}\xspace}
\newcommand{\unit}{\mylink{\finitetypecolor}{unit}{\unitnolink}\xspace}
\newcommand{\unite}{\ensuremath{\mylink{\finitetypecolor}{unit}{\star}}\xspace}
\newcommand{\finel}[2]{\mylink{\finitetypecolor}{emptytype}{\hookrightarrow}_{#2}#1}
\newcommand{\enumsym}{\ensuremath{\mylink{\finitetypecolor}{enum}{\textsf{enum}}}\xspace}
\newcommand{\enum}[1]{\enumsym\mylink{\finitetypecolor}{enum}{\left( \vphantom{#1} \right. } #1 \mylink{\finitetypecolor}{enum}{\left. \vphantom{#1}\right) }}
\newcommand\enumusym{\ensuremath{\mylink{\finitetypecolor}{enum}{\textsf{case}}}\xspace}
\newcommand{\enume}[2]{\enumusym_{#2}\mylink{\finitetypecolor}{enum}{\left( \vphantom{#1} \right.}#1\mylink{\finitetypecolor}{enum}{\left. \vphantom{#1} \right) }}
\newcommand{\enumu}[1]{\enumusym\mylink{\finitetypecolor}{enum}{\left( \vphantom{#1}\right.}#1\mylink{\finitetypecolor}{enum}{\left. \vphantom{#1}\right)}}
\newcommand\numtp[1]{\mylink{\finitetypecolor}{unit}{\textbf{#1}}}
\newcommand\numel[1]{\mylink{\finitetypecolor}{unit}{\underline{#1}}}

% W-Types------------------------------------------------------------------------------------------------------------------------------
%wtype,wsup,wrec

\newcommand{\wtypecolor}{Melon}

\DeclareMathOperator*{\wtypesymm}{\mylink{\wtypecolor}{wtype}{\textsf{\larger W}}}
\newcommand{\wtypesym}{\ensuremath{\wtypesymm}}
\newcommand{\wtype}[2]{\wtypesym_{#1} #2}
\newcommand{\wsupsym}{\mylink{\wtypecolor}{wsup}{\ensuremath{\sup}}\xspace}
\newcommand{\wsup}[3]{\wsupsym_{#1}\mylink{\wtypecolor}{wsup}{\left\{ \vphantom{#2 #3}\right.} #2 \mylink{\wtypecolor}{wsup}{\Longrightarrow} #3\mylink{\wtypecolor}{wsup}{\left. \vphantom{#2 #3}\right\}}}

\newcommand{\wsupu}[4]{
\wsupsym_{#1}\mylink{\wtypecolor}{wsup}{\left\{\vphantom{#2 #3 #4}\right. } #2 \mylink{\wtypecolor}{wsup}{\Longrightarrow}_{#4} #3\mylink{\wtypecolor}{wsup}{\left.\vphantom{#2 #3 #4}\right\}} }

\newcommand{\wrecsym}{\ensuremath{\mylink{\wtypecolor}{wrec}{\textsf{rec}}}\xspace}
\newcommand{\wrec}[7]{\wrecsym\mylink{\wtypecolor}{wrec}{ \left( \vphantom{#2} \right.}#2\mylink{\wtypecolor}{wrec}{\left. \vphantom{#2}\right)}\mylink{\wtypecolor}{wrec}{\left\{ \left( \vphantom{#4 #5 #6} \right. \vphantom{#7} \right. } #4,#5,#6\mylink{\wtypecolor}{wrec}{\left. \vphantom{#4 #5 #6} \right)} \mylink{\wtypecolor}{wrec}{\Longrightarrow}_{#3} #7 \mylink{\wtypecolor}{wrec}{\left. \vphantom{#4 #5 #6 #7} \right\} } }

% Hierarchy-----------------------------------------------------------------------------------------------------------------------------
%univ

\definecolor{darkred}{RGB}{100,0,0}
\newcommand{\univcolor}{darkred}

\newcommand\univ[1]{\mylink{\univcolor}{univ}{{\mathcal U}}_{#1}}

% Equality Types----------------------------------------------------------------------------------------------------------------------
%eq

\newcommand{\eqcolor}{Tan}

\newcommand\eqtp[3]{#1\mylink{\eqcolor}{eqtp}{\,\doteq\,}_{#3}#2}
\newcommand\refl[1]{\mylink{\eqcolor}{eqtp}{\mathtt{refl}}_{#1}}
\newcommand\eqindsym{\mylink{\eqcolor}{eqtp}{\ensuremath{\mathtt{ind}}}\xspace}
\newcommand\eqind[6]{\eqindsym_{#1} \mylink{\eqcolor}{eqtp}{\left\{ \vphantom{#6,#2,\ _{#5},#3,#4} \right.} #2,#3,#4 \mylink{\eqcolor}{eqtp}{\Longrightarrow}_{#5} #6 \mylink{\eqcolor}{eqtp}{\left. \vphantom{#6,#2,\ _{#5},#3,#4}\right\}} }

% Subypes----------------------------------------------------------------------------------------------------------------------
%stp,inttp,predsubtp

\newcommand{\stpcolor}{CadetBlue}

\newcommand\stpsym{\ensuremath{\mylink{\stpcolor}{stp}{<\ast}}}
\newcommand\stp[2]{\ensuremath{#1 \,\stpsym\, #2 }}
\newcommand\inttpsym{\ensuremath{\mylink{\stpcolor}{inttp}{\sqcap}}}
\newcommand\inttp[2]{\ensuremath{#1 \,\inttpsym\, #2}}
\newcommand\predsubtp[3]{ 
	\mylink{\stpcolor}{predsubtp}{\left\langle \vphantom{ #1 #2 #3} \right.}
		#1 : #2 \; \mylink{\stpcolor}{predsubtp}{\left\mid \vphantom{ #1 #2 #3} \right.} \;
		#3 \mylink{\stpcolor}{predsubtp}{\left. \vphantom{ #1 #2 #3} \right\rangle}
}
\newcommand\predofsym{\mylink{\stpcolor}{predsubtp}{\ensuremath{\mathtt{predOf}}}}
\newcommand\predof[1]{\mylink{\stpcolor}{predsubtp}{\predofsym\left( \vphantom{#1} \right. } #1
\mylink{\stpcolor}{predsubtp}{\left. \vphantom{#1} \right)}
}

% Records (TODO) -------------------------------------------------------------------------------------------------------------------
%records, models, subsume

\newcommand{\recordcolor}{RubineRed}

\newcommand{\leftrectype}[1]{\mylink{\recordcolor}{records}{\left\lsem\vphantom{#1}\right.}}
\newcommand{\rightrectype}[1]{\mylink{\recordcolor}{records}{\left.\vphantom{#1}\right\rsem}}
\newcommand{\rectype}[1]{\leftrectype{#1} #1 \rightrectype{#1}}
%\newcommand{\leftrecexp}{\mylink{\recordcolor}{records}{\left\langlebar}}
%\newcommand{\rightrecexp}{\mylink{\recordcolor}{records}{\right\ranglebar}}
\newcommand{\recsym}{\,\mylink{\recordcolor}{records}{\mathrlap{\succ}}-\,}
\newcommand{\recexp}[1]{\mylink{\recordcolor}{records}{\left\langlebar\vphantom{#1}\right.} #1 \mylink{\recordcolor}{records}{\left.\vphantom{#1}\right\ranglebar}}
\newcommand{\recexpi}[2]{\recexp{ #1 \recsym #2}}
\newcommand{\self}{\mylink{\recordcolor}{records}{\mathtt{self}}}
\newcommand{\recproj}[2]{#1\mylink{\recordcolor}{records}{.}#2}

\newcommand{\modelsofF}{\mylink{\recordcolor}{models}{\mathtt{Mod}}}
\newcommand{\modelsof}[1]{\modelsofF\mylink{\recordcolor}{models}{\left(\vphantom{#1}\right.}#1\mylink{\recordcolor}{models}{\left.\vphantom{#1}\right)}}
\newcommand{\subsume}[2]{#1 \mylink{\recordcolor}{subsume}{\harr} #2}

\newcommand{\recmergesym}{\mylink{\recordcolor}{merge}{+}}
\newcommand{\recmerge}[2]{#1\recmergesym #2}
\newcommand{\contmergesym}{\mylink{\recordcolor}{merge}{\mp}}
\newcommand{\contmerge}[2]{#1\contmergesym #2}

% Others----------------------------------------------------------------------------------------------------------------------------------

\newcommand{\nat}{\ensuremath{\textsc{Nat}}\xspace}
\newcommand{\sub}[2]{\mylink{Bittersweet}{subst}{\left[ \vphantom{#1 /_{#2}} \right. } #1 \mylink{Bittersweet}{subst}{/}_{#2}\mylink{Bittersweet}{subst}{\left. \vphantom{#1 /_{#2}} \right]}}

% Old? -------------------------------------------------------------------------------------------------------------------------------------

\newcommand{\Log}{\cn{Log}}
\newcommand{\keyword}[1]{\ensuremath{\operatorname{\mathbf{#1}}}}
\newcommand{\imp}{\hookrightarrow}
\newcommand{\hvee}{\sqcup}
\newcommand{\hwedge}{\sqcap}
\newcommand{\set}[1]{\left\lbrace #1 \right\rbrace}
\newcommand{\abs}[1]{\left| #1 \right|}
\newcommand{\mm}{\mathfrak{mm}}
\newcommand{\pid}{\textsf{PID}}
\newcommand{\pfa}{\textsf{PFA}}
\newcommand{\gch}{\textsf{GCH}}
\newcommand{\PI}{\mathcal P_{\mathcal I}}
\newcommand{\ipc}[3]{#2\stackrel{#1}{\harr}#3}
\newcommand{\ipce}[2]{#1\stackrel{\ast}{\harr}#2}
\newcommand{\icl}[1]{\incl{#1}}
\newcommand{\dom}[1]{\mathtt{dom}(#1)}
\newcommand{\flt}[1]{#1^\flat}
\newcommand{\Exp}[1]{\mathtt{Obj}(#1)}

\newcommand{\flti}[1]{#1^{\flat i}}
\newcommand{\Expi}[1]{\mathtt{Obj}^i(#1)}
\newcommand{\smt}[1]{\sqrbrc{#1}}
\newcommand{\thy}{\mathtt{thy}}
\newcommand{\Thy}{\mathtt{Thy}}
\newcommand{\mo}[2]{\mathtt{mor}(#1,#2)}
\newcommand{\Mo}[2]{\mathtt{Mor}(#1,#2)}
\newcommand{\IMo}[2]{\mathtt{Mor^i}(#1,#2)}
\newcommand{\strc}[2]{\keyword{struct}{\;#1:#2}}
\newcommand{\fltd}[2]{#2^{\flat}}
\newcommand{\cls}[1]{\ov{#1}}


\renewcommand{\bnf}[1]{{\color{red}#1}}

%\newcommand{\recsub}[2]{[#1 / #2]}
\def\defemph{\textbf}

\newcommand{\SL}{\mathtt{Semilattice}}
\newcommand{\SLO}{\mathtt{SemilatticeOrder}}
\newcommand{\csl}{\mathtt{interSL}}
\newcommand{\cslo}{\mathtt{interSLO}}

\def\cB{\mathcal{B}}
\def\cL{\mathcal{L}}


\newcommand{\gray}[1]{{\color{gray}#1}}
\newcommand{\red}[1]{{\color{red}#1}}


\def\Domain#1{\text{Dom}(#1)}
\def\conc{::}